\section{Discussion and limitations}
\label{sec:discussion}

\subsection{What the finite-rank comparison establishes}

The calculations answer a question that cannot be posed through ordinary
level statistics.  When a parent Hamiltonian has a rank-$D$ eigenspace pinned
to one energy, neither its zero spacings nor its spectral form factor
distinguishes a rigid projector from one that rotates strongly under local
changes of couplings.  The non-Abelian quantum geometric tensor resolves this
difference because it measures virtual excursions from the protected fiber
into the gapped complement.  The relevant randomness is therefore randomness
of the embedding $P(\boldsymbol\lambda)\subset\mathcal H$, not randomness of
the repeated eigenvalue $E_0$.

Across clustered Hall parents and cohomological kernels, the response is
nontrivial even though the protected energies remain equal at the accepted
base point or along an exact transported family.  Several finite-rank
curvature spectra display Jacobi-like local eigenvalue repulsion.  This is the
part of the result that most closely resembles an ordinary random-matrix
criterion: nearby curvature eigenvalues avoid one another after the
finite-$D$ density is treated correctly.  It is also the narrowest part of the
claim.  The comparison does not show that the full response matrix is drawn
from an invariant Jacobi ensemble.  Its channel covariance, pseudocovariance,
and exact/coexact or stabilizer branch structure remain visible at larger
scales.

We use \emph{geometric chaos} for this combination of noncommuting projector
response, stochastic local curvature correlations, and a nontrivial
complement channel.  The phrase does not assert that the parent Hamiltonian
thermalizes.  Indeed, the protected sector has no internal energy differences
on which a conventional ETH ansatz could act.  What is measured is the
sensitivity and correlation structure of the degenerate subspace under a
declared tangent ensemble.  This meaning is consistent with the use of
eigenstate deformation as a chaos-sensitive observable \cite{pandey2020} and
with the proposal to use Berry curvature inside BPS multiplets
\cite{chen2026}, while remaining agnostic about long-time relaxation.

\subsection{Protection, mixing, and the meaning of a tangent}

The separation between $S_a$ and $X_a$ is more than a bookkeeping device.
$S_a$ asks whether a perturbation resolves energies inside the fiber;
$X_a=Q(\partial_aP)P$ asks whether it changes the fiber's embedding.  An
exactly transported parent can have $S_a=0$ and large $X_a$, whereas a scalar
rescaling can have both responses trivial.  Conversely, an arbitrary local
potential can split the fiber and rotate its spectral cluster at the same
time.  Calling all three operations \emph{deformations of the parent} would
erase the physical distinction on which the geometric test rests.

Accordingly, finite-path claims in this work come only from
exactness-preserving transport.  For the Laughlin and Moore--Read generator
panels, the constraint map is transported unitarily, the kernel rank remains
fixed, and the differentiated identity gives $S_a=0$.  Closed paths can then
support a Wilczek--Zee holonomy \cite{wilczekzee1984}; its non-Abelian part
contains information not fixed by the first Chern number.  In a gapped
topological phase, quasiadiabatic continuation supplies the broader setting
for such projector motion \cite{hastingswen2005}, although our exact parent
paths are more restrictive than a generic phase-preserving deformation.

The Laughlin local-potential panels are fixed-parent base-point probes.
Their $Q(\partial_aH)P$ matrices are well defined because the rank-$D$
spectral cluster is isolated at the parent, but a generic finite potential
does not preserve exact equality of its eigenvalues.  These panels can test
local response statistics at the exact point; they cannot support a claim
about a globally degenerate bundle away from it.  This boundary is especially
important when interpreting a large curvature response: sensitivity of an
isolated cluster is not by itself protection of a degenerate multiplet.

\subsection{Local Jacobi behavior and global memory}

A finite Jacobi process is an appropriate reference for eigenvalues generated
by the relative orientation of random subspaces \cite{collins2005}.  Agreement
with its local repulsion law tests correlations within a single whitened
curvature matrix.  It does not test every index of $X_a$, nor does it remove
the covariance between tangent directions.  The analogy with standard
spectral chaos must therefore stop before one equates a local point process
with a complete matrix distribution.

This separation explains why Jacobi-like local eigenvalue repulsion can
coexist with global covariance memory.  A clustered parent restricts which
complement states are reached by a local generator.  A supercharge resolves
the response into exact and coexact branches.  A stabilizer algebra confines
the response to a small commutant-determined sector.  These restrictions
survive covariance whitening unless the whitening includes the complete entry
covariance and pseudocovariance with the correct independent sampling unit.
They can also generate fourth-order structure that is invisible in a
nearest-neighbor spacing histogram.

For that reason, the four-point analysis uses all three complex Wick pairings
and treats the full fixed-Hamiltonian tangent panel as one statistical object.
Matrix entries and pairs of labels within that panel are not independent
realizations.  A nonzero raw connected estimate obtained by resampling them as
if they were independent would measure pseudoreplication as readily as
non-Gaussian physics. The final evidence registry therefore uses the exactly enumerated population mean of each registered finite ensemble, a channel whitener fitted on a frozen 12-panel training split, and the complete three-pairing statistic on 12 held-out panels. The primary inference treats those held-out panels as exchangeable conditional on the fixed base Hamiltonian, uses a Student-$t_{11}$ reference, and applies a Bonferroni familywise correction over the five registered Moore--Read and lattice-SUSY cases. The reverse 12/12 split is descriptive and is not pooled with the primary inference.

Only the registered finite-rank Moore--Read $N=6$ case passes this corrected gate: its directional estimate is $\MooreReadNSixFullRFourDirectionalEstimate$, and its familywise one-sided lower bound is $\MooreReadNSixFullRFourFamilywiseLower>0$. The Moore--Read $N=4$ lower bound remains negative, as do those for lattice-SUSY ranks $D=2,4,8$. No positive complete-cumulant statement is made for those nonpassing cases, and nonpassing is not evidence for Gaussianity. The Laughlin and random-SYK production data have not been evaluated with this corrected statistic. The $N=6$ result is therefore evidence for a population-centered fourth-order residual in one finite-rank, fixed-base tangent ensemble; it is not evidence for an independent-Hamiltonian ensemble or a common large-rank law.

\subsection{What is mechanism dependent}

The cross-mechanism comparison contains both positive and negative cases.
Two- and three-body clustering produce large common kernels whose response is
set by forbidden coincidence channels.  Random-supercharge cohomology gives a
dense coupling-space ensemble with exact and coexact Hodge branches.  The
local lattice supercharge preserves cohomology but has a different locality
and a much smaller sequence of protected ranks.  These systems need not share
a covariance tensor, and the present data do not support pooling their
uncertainties into one scaling curve.

The X-cube model is the decisive structured control.  Changing coefficients
of fixed commuting stabilizers leaves the eigenprojector unchanged, so the
coefficient-space curvature is zero.  Under a specified local unitary
transport the projector moves and its curvature can be nonzero, but the
curvature eigenvalue is deterministic and its connected variance vanishes.
This example separates three propositions that are sometimes conflated:
macroscopic exact degeneracy, nonzero geometric response, and stochastic
geometric chaos.  Neither of the first two implies the third.

The continuum and Kapit--Mueller Laughlin calculations give a second kind of
control.  They share two-body clustering but differ microscopically.  Agreement
of a local statistic across them is evidence that the effect is not solely a
lattice artifact; disagreement of a global covariance contraction is allowed
because the tangent realization and complement spectrum differ.  Parent
dependence is consequently part of the result rather than an error bar to be
averaged away.

\subsection{Relation to Geometric ETH}

For the finite systems studied here, \emph{Geometric ETH} is an organizing name
for a hierarchy of tests, not a theorem.  The weakest level is informative
geometry under exact degeneracy.  A stronger, domain-limited level adds local
random-matrix-like correlations for a specified protection mechanism and
tangent class.  An asymptotic statement would require a sequence with growing
protected rank, a fixed normalization and operator class, controlled external
gaps, converged local statistics, and decay or convergence of the
covariance-whitened connected tensors.  None of these requirements can be
replaced by increasing the ambient Hilbert-space dimension alone.

Our result is finite-rank, cross-mechanism evidence.  It is not an asymptotic
Geometric ETH theorem and not a universal Geometric ETH theorem.  In
particular, an independent model/operator class has not been established:
the available sequences change either the protection algebra, the tangent
panel, or both.  We also do not demonstrate thermalization, a Lyapunov
exponent, an out-of-time-order growth regime, or equivalence to the usual ETH
ansatz for few-body observables.  Those are distinct dynamical propositions.

These limitations indicate what would raise the claim.  Within one mechanism,
one must extend a fixed operator prescription through a substantially larger
rank sequence and preserve a positive external gap.  Independent Hamiltonian
or disorder realizations must replace fixed-base panels wherever an ensemble
claim requires them.  Complete covariance and pseudocovariance must be
estimated with uncertainty at that realization level.  Only after these
conditions are met would it be meaningful to test a common large-$D$ law.
Until then, mechanism-complete sampling is more informative than adding an
unbounded list of models with incompatible tangent definitions.

\subsection{Conclusion}
\label{sec:conclusion}

Exactly degenerate energies do not make quantum chaos unobservable; they move
the question from eigenvalue spacings to the geometry of the protected
projector.  The response $X_a=Q(\partial_aP)P$ remains finite when the
intrafiber splitting channel vanishes, and its non-Abelian curvature separates
rigid protected fibers from stochastically mixing ones.  Across clustered Hall
parents, random and local cohomological models, and the X-cube control, we find
finite-rank, cross-mechanism evidence for a hierarchy in which locally
Jacobi-like curvature can retain parent- and operator-dependent covariance at
larger scales.

The hierarchy is the main conclusion.  Local repulsion is not full Gaussian
closure; nonzero curvature is not stochasticity; base-point sensitivity is not
finite-path protection; and exact degeneracy is not thermalization.  The
negative stabilizer result is as important as the stochastic cases because it
shows that exact protection does not automatically generate a random geometry.
The positive statements are therefore conditional on the declared mechanism,
tangent class, protected rank, and statistical unit.

The present evidence supports a finite-size, domain-limited use of the term Geometric ETH. It leaves the asymptotic problem open. Establishing that stronger statement will require growing-rank sequences with an unchanged operator prescription and independent-ensemble covariance closure. Results from the corrected full-cumulant gate are restricted to Moore--Read $N=6$; Moore--Read $N=4$ and lattice-SUSY $D=2,4,8$ do not pass, while Laughlin and random SYK remain untested by that statistic. This finite-rank result does not establish an asymptotic or universal Geometric ETH. Within these boundaries, non-Abelian quantum geometry provides information precisely where the degenerate spectrum does not.
